\chapter{Literature Review}
\label{ch:literature}

\section{Document Understanding}
\label{sec:doc_understanding}

Document understanding is the task of converting unstructured document images
into structured information. It supports information extraction and automated
document-processing workflows. Earlier systems often combined \gls{ocr} with
hand-crafted rules, regular expressions, keyword dictionaries, and layout
templates. Representative systems used structural templates or user-provided
examples to locate fields in administrative and business
documents~\cite{rusinol2013field,schuster2013intellix}. Such systems often
require manual configuration or template updates when a document layout
changes~\cite{rusinol2013field,schuster2013intellix,layoutlm}.

\Gls{cnn} models were applied to document recognition and classification to
learn visual features from document images~\cite{lecun1998gradient,rvlcdip}.
\Gls{lstm} networks provided a recurrent architecture for sequence
modelling~\cite{hochreiter1997lstm}. The Transformer replaced recurrence with
self-attention~\cite{vaswani2017attention}. Later layout-aware document models
combined text, page geometry, and image features.

\section{Multimodal Pre-training for Documents}
\label{sec:multimodal}

Multimodal document pre-training learns joint representations of text, page
geometry, and image content. These representations can then be adapted to
downstream document tasks~\cite{layoutlm,layoutlmv2,layoutlmv3}. The LayoutLM
family is one prominent model family in this area.

LayoutLM~\cite{layoutlm} extends Bidirectional Encoder Representations from
Transformers (BERT)~\cite{transformers} with
two-dimensional bounding-box embeddings for document tokens. It was
pre-trained on the Illinois Institute of Technology Complex Document
Information Processing (IIT-CDIP) collection~\cite{iitcdip}. The original
study evaluated it on Form
Understanding in Noisy Scanned Documents (FUNSD)~\cite{funsd}, the Scanned
Receipt OCR and Information Extraction (SROIE) benchmark~\cite{sroie}, and the
Ryerson Vision Lab Complex Document Information Processing (RVL-CDIP)
dataset~\cite{rvlcdip}. The results show that explicit page geometry can
improve document tasks when it is combined with text.

LayoutLMv2~\cite{layoutlmv2} adds document-image features from a \gls{cnn}
backbone. It also uses spatial-aware self-attention based on relative token
geometry. Its pre-training objectives jointly use text, layout, and image
information. The original study reports improvements across its
document-understanding benchmarks.

LayoutLMv3~\cite{layoutlmv3} replaces the convolutional image encoder with
image patches, following the Vision Transformer design~\cite{vit}. It uses
masked language modelling, masked image modelling, and word-patch alignment
during pre-training. At release, it reported state-of-the-art results on
FUNSD~\cite{funsd}, CORD (Consolidated Receipt Dataset for Post-OCR
Parsing)~\cite{cord}, and Document Visual Question Answering
(DocVQA)~\cite{docvqa}. LayoutLMv3 is the pretrained encoder used by the
discriminative system in this thesis.

\section{Key-Value Pair Extraction}
\label{sec:kvp_extraction}

Key-value pair extraction is the main task in this thesis. Given a document
page, a system must identify key and value entities, localize them on the page,
and determine which keys and values form pairs. Different approaches divide
these subtasks in different ways.

One common formulation treats entity detection as token classification. The
\gls{bio} scheme marks whether a token starts an entity, continues an entity,
or lies outside an entity~\cite{ramshaw1995text}. LayoutLM-based systems use
token-level sequence labelling for form and receipt information
extraction~\cite{layoutlm,layoutlmv2,layoutlmv3}. FUNSD also provides relations
between its annotated semantic entities~\cite{funsd}. In contrast, the
discriminative system in this thesis assigns \texttt{KEY}, \texttt{VALUE}, or
\texttt{O} to each token, where \texttt{O} identifies tokens that belong to
neither entity class. A separate relation model then links detected keys and
values.

Generative approaches serialize a document and ask an autoregressive language
model to produce a structured extraction output~\cite{seq2seq,lmdx,kvp10k}.
This formulation supports open-vocabulary output, but the generated output must
be parsed and grounded to document locations.

The \gls{lmdx} method represents OCR segments as text followed by quantized
coordinate tokens. Its prompt also contains a task description and a target
schema. A decoder maps generated text and location tokens back to document
entities and rejects outputs that cannot be grounded in the
document~\cite{lmdx}.

The published KVP10k Mistral-7B generative baseline adapts this idea to
open-vocabulary key-value pair extraction~\cite{kvp10k,mistral}. It orders OCR
text from top to bottom and left to right and appends the bounding-box
coordinates \texttt{[x1,y1,x2,y2]} to each text line. Mistral-7B is then
fine-tuned to generate the benchmark output. The output is a sequence of
dictionaries. Each entry has a type (\texttt{kvp}, \texttt{unkeyed}, or
\texttt{unvalued}) and the available key or value text and bounding boxes.

This thesis implements a reconstructed Mistral-7B reference system that follows
this high-level approach. It is not an exact reproduction of the published
system. Later chapters describe the local prompt, Python list-of-lists
serializer, fine-tuning procedure, output parser, and evaluator integration.

Parameter-efficient fine-tuning adapts a pretrained model by updating a small
set of additional parameters. \Gls{lora} inserts trainable low-rank matrices
into selected model layers while the base weights remain fixed~\cite{lora}.
\Gls{qlora} combines low-rank adaptation with a quantized base model to reduce
memory use during fine-tuning~\cite{qlora}. The reconstructed Mistral-7B
reference system in this thesis uses QLoRA.

Graph-based methods represent detected text regions as nodes and spatial or
semantic relations as edges. PICK (Processing Key Information Extraction from
Documents using Improved Graph Learning-Convolutional Networks)~\cite{pick}
builds a document graph and learns contextual node representations for key
information extraction. Graph methods can model relations between distant text
regions, but they require a graph definition and an explicit relation
structure.

\subsection{Entity Linking Approaches}
\label{subsec:entity_linking}

Another family of methods separates entity detection from entity linking. The
first component identifies key and value spans. The second component predicts
which detected key and value belong to one pair~\cite{spade,bros}. This
decomposition gives each subtask a separate model and matches the general
classifier-linker design used in this thesis.

SPADE (Spatial Dependency Parsing for Semi-Structured Document Information
Extraction)~\cite{spade} treats document information extraction as dependency
parsing over OCR tokens. It predicts directed token relations that assign
fields and connect tokens into a structured output graph. This design supports
two-dimensional layouts and hierarchical outputs without a fixed
one-dimensional reading order. SPADE provides method context for relation-based
extraction. The contiguous \texttt{KEY}/\texttt{VALUE} span grouping used by
the span-level discriminative variants in this thesis is a separate
implementation decision (Section~\ref{sec:v2-span-linking}).

BROS (BERT Relying On Spatiality)~\cite{bros} combines text with relative
two-dimensional positions between text-block bounding-box vertices. It adds
area-masked language modelling during pre-training. For entity linking, BROS
uses an adapted SPADE decoder that predicts token relations. The study reports
results on document benchmarks including FUNSD~\cite{funsd} and
CORD~\cite{cord}. BROS provides context for relative spatial encoding and
relation decoding, but it does not define the span grouping or candidate
reduction used in this thesis (Section~\ref{sec:v2-span-linking}).

Relation models require a score for each candidate pair. Biaffine scoring uses
a learned bilinear interaction with affine terms and is used in structured
prediction~\cite{dozat2017biaffine}. Scaled dot-product attention compares
projected query and key representations and divides their dot product by
$\sqrt{d_k}$, where $d_k$ is the query and key
dimension~\cite{vaswani2017attention}. These methods provide the technical
context for the change from the initial biaffine linker to the projected
scaled-dot-product relation scorer in this thesis.

Recent document models add task-specific relation components. GeoLayoutLM
pre-trains geometric relation objectives and relation heads~\cite{geolayoutlm}.
KVPFormer treats detected keys as questions and predicts value entities with a
coarse-to-fine decoder and a spatial compatibility bias~\cite{kvpformer}.
LayoutPointer uses relative two-dimensional context, adaptive spatial
constraints, and pointer-based relation prediction for multi-line
entities~\cite{layoutpointer}. These methods show different ways to model
spatial entity association explicitly.

\section{Information Extraction from Forms}
\label{sec:form_extraction}

Forms and receipts are common document information extraction domains. Their
datasets support research on entity recognition, relation extraction, and
semantic parsing~\cite{funsd,cord,sroie}.

FUNSD contains 199 scanned forms with word-level text, semantic labels, and
entity-linking annotations~\cite{funsd}. CORD contains more than 11,000
collected receipts, and its released benchmark contains 1,000 annotated samples
for receipt parsing~\cite{cord}. SROIE provides a receipt OCR and information
extraction benchmark~\cite{sroie}. These datasets focus on forms or receipts,
while KVP10k contains a broader set of business-document
layouts~\cite{kvp10k}.

\section{Evaluation Metrics}
\label{sec:metrics_review}

The KVP10k benchmark evaluates both text and location
accuracy~\cite{kvp10k}. Text matching uses \gls{ned}, which is based on edit
distance~\cite{levenshtein}. Location matching uses
\gls{iou}~\cite{iou}. For the key-value pair detection task, the paper defines
a text hit as NED below $0.2$ and a location hit as IoU above $0.3$. It reports
text-only, location-only, and text+location scores and macro-averages precision
and recall per page~\cite{kvp10k}.

An inspection of the released evaluator~\cite{kvp10k_repository} shows that it
accepts an IoU value of exactly $0.3$, because its implemented boundary is inclusive ($\geq 0.3$).
This thesis follows the released evaluator for all reported comparisons.
Section~\ref{sec:evaluation-protocol} gives the formal definitions and records
this paper-code difference.

\section{Related Datasets}
\label{sec:related_datasets}

Document-understanding datasets cover different tasks, document types, and
output structures.

KVP10k~\cite{kvp10k} is a benchmark for open-vocabulary key-value pair
extraction from business documents. The paper reports 10{,}707 pages. In the
released annotations used in this thesis, grouping rows by the page identifier
\texttt{hash\_name} gives 9{,}656 unique training pages and 1{,}051 unique test
pages, as detailed in Section~\ref{sec:kvp10k}. Its annotations provide text
regions, page-relative coordinates, and links from values to their keys. The
benchmark includes regular key-value pairs, unkeyed values, and unvalued keys.
It also provides text-only, location-only, and text+location evaluation modes.
This task definition and dataset scale make KVP10k the main dataset for this
thesis.

DocVQA asks systems to answer natural-language questions about document
images~\cite{docvqa}. RVL-CDIP contains 400,000 document images in 16 classes
for document-image classification~\cite{rvlcdip}. XFUND is a multilingual
form-understanding benchmark for seven languages other than
English~\cite{xfund}. These datasets are not used in the experiments. They
provide context for question answering, document classification, and
multilingual form understanding.

\section{Scope of Experimental Design}
\label{sec:experimental_scope}

The experiments use two system families: a LayoutLMv3-based discriminative
classifier-linker pipeline and a reconstructed Mistral-7B reference system. The
study does not reimplement SPADE, BROS, GeoLayoutLM, KVPFormer, or
LayoutPointer. These methods therefore provide literature context and possible
future baselines. The direct comparison between the reconstructed Mistral-7B
reference system and the final LayoutLMv3-based span-level discriminative model
(V4) uses the same prepared KVP10k test subset and the same released evaluator.

\section{Summary and Thesis Objectives}
\label{sec:lit_summary}

The review motivates the following thesis objectives:

\begin{enumerate}
    \item \textbf{Controlled KVP10k evaluation}: This thesis evaluates the
    final LayoutLMv3-based span-level discriminative model (V4) and the
    reconstructed Mistral-7B reference system with the released KVP10k
    evaluator. Their direct comparison uses the same prepared test subset.

    \item \textbf{Separate text and location analysis}: The KVP10k metrics
    separate text-only, location-only, and text+location matching. This thesis
    uses these modes to compare performance under text-only, location-only, and
    text+location matching criteria.

    \item \textbf{Development diagnostics}: Token-level V1 and span-level V2
    are development variants. Their validation diagnostics examine candidate
    construction, bounding-box filtering, and relation supervision. V4 is the
    final span-level system used for the main discriminative evaluation. These
    diagnostics do not form a controlled causal comparison between all
    variants.

    \item \textbf{Geometry-based error analysis}: This thesis examines
    performance across exploratory page clusters based on annotation-box
    count, size, and spatial spread.
\end{enumerate}

LayoutLMv3 supplies the pretrained encoder, and the general scaled-dot-product
scoring principle is prior work~\cite{layoutlmv3,vaswani2017attention}. The
task-specific classifier and linker modules, span-grouping path, training
pipeline, evaluator integration, and diagnostic procedures were implemented
for this thesis. SPADE, BROS, and the newer relation models provide method
context, but they were not reimplemented. V1 uses token-level candidates, V2
is an early span-level development variant, and V4 is the final span-level
model.
